%% ============================================================
%% Title Page — PAR Submission (separate from blind manuscript)
%% ============================================================

\documentclass[12pt]{article}
\usepackage[margin=1in]{geometry}
\usepackage{setspace}
\doublespacing
\usepackage{times}
\usepackage{hyperref}
\hypersetup{colorlinks=true, linkcolor=black, citecolor=black, urlcolor=black}

\date{}
\author{}

\begin{document}

\begin{center}

\vspace*{1.5cm}

{\LARGE\textbf{The Researcher-Agent Problem: A Principal-Hierarchy Framework for
AI-Assisted Knowledge Production in Public Administration}}

\vspace{1.5cm}

\textbf{Edmund Poku Adu, PhD}\\[4pt]
Assistant Professor\\
Department of Government, Law, and Policy\\
Arkansas State University\\
P.O. Box 1750, State University, AR 72467\\
\href{mailto:eadu@astate.edu}{eadu@astate.edu}

\vspace{2cm}

\textbf{Corresponding author:} Edmund Poku Adu, PhD\\
\href{mailto:eadu@astate.edu}{eadu@astate.edu}

\vspace{2cm}

\textbf{Word count:} approximately 6,100 words (excluding abstract, tables, and references)

\vspace{1cm}

\textbf{Acknowledgments:} The author thanks Dr.\ Valencia Prentice (Cleveland State
University) for an extended conversation that served as the motivating spark for this
article.

\vspace{1cm}

\textbf{Data availability:} All session transcripts, workflow state files, and output
artifacts supporting the case study findings are deposited at
\url{https://github.com/epokuadu/PA-Agent} and are publicly available for replication.

\vspace{1cm}

\textbf{AI disclosure:} This manuscript was produced using the PA-Agent agentic research
workflow, which is the subject of the case study presented herein. All researcher-agent
interactions are fully documented in the replication repository. The researcher-principal
retains full responsibility for all claims, interpretations, and conclusions.

\end{center}

\end{document}
